\documentclass[11pt,a4paper]{article}
\usepackage[utf8]{inputenc}
\usepackage[T1]{fontenc}
\usepackage[ngerman]{babel}
\usepackage{helvet}
\renewcommand{\familydefault}{\sfdefault}
\usepackage[left=2cm, right=2cm, top=2cm, bottom=2cm]{geometry}
\usepackage{amsmath,amssymb}
\usepackage{array}
\usepackage{booktabs}
\usepackage{tabularx}
\usepackage{tikz}
\usepackage{pgfplots}
\pgfplotsset{compat=1.18}
\usepackage{fancyhdr}
\usepackage{siunitx}
\sisetup{locale = DE}
\usepackage{xcolor}
\usepackage{marginnote}
\reversemarginpar

\pagestyle{fancy}
\fancyhf{}
\fancyhead[L]{\small Physik Klasse 10E}
\fancyhead[R]{\small Probeklausur Kinematik}
\fancyfoot[C]{\small Seite \thepage}
\renewcommand{\headrulewidth}{0.4pt}

\setlength{\parindent}{0pt}
\setlength{\fboxsep}{8pt}

% Farben
\definecolor{afbgray}{RGB}{140,140,140}
\definecolor{zeitgray}{RGB}{120,120,120}

% BE + Zeit + AFB rechtsbündig
\newcommand{\BEZEIT}[3]{%
\hfill{\small\color{zeitgray}(#1\,Min, AFB\,#2)}\quad\textbf{(\_\_/#3\,BE)}%
}

\begin{document}

% === KOPFBEREICH ===
\noindent\fbox{\parbox{\dimexpr\textwidth-2\fboxsep-2\fboxrule}{
\centering
{\Large\bfseries Probeklausur: Kinematik}\\[0.3cm]
{\normalsize Klasse 10E \quad|\quad 90 Minuten \quad|\quad 48 BE}\\[0.2cm]
{\small\textit{Diese Übungsklausur dient der Vorbereitung auf die Klausur.}}
}}

\vspace{0.4cm}

\begin{tabular}{@{}p{6cm}p{8cm}@{}}
Name: & \hrulefill \\[0.3cm]
Datum: & \hrulefill \\
\end{tabular}

\vspace{0.3cm}

% === NOTENSPIEGEL (15-NP Sek II MV) ===
\begin{center}
\small
\begin{tabular}{|c|c|c|c|c|c|c|c|c|c|c|c|c|c|c|c|}
\hline
\textbf{NP} & 15 & 14 & 13 & 12 & 11 & 10 & 9 & 8 & 7 & 6 & 5 & 4 & 3 & 2 & 1 \\
\hline
\textbf{ab BE} & 46 & 43,5 & 41 & 38,5 & 36 & 34 & 31,5 & 29 & 26,5 & 24 & 22 & 19,5 & 16 & 13 & 10 \\
\hline
\textbf{ab \%} & 95 & 90 & 85 & 80 & 75 & 70 & 65 & 60 & 55 & 50 & 45 & 40 & 33 & 27 & 20 \\
\hline
\end{tabular}
\end{center}

\vspace{0.3cm}

\begin{center}
\fbox{\parbox{0.95\textwidth}{
\small
\textbf{Hinweise:}
\begin{itemize}
    \item Erlaubte Hilfsmittel: Tafelwerk, Geodreieck, Lineal, Bleistift
    \item \textbf{Taschenrechner ist NICHT zugelassen}
    \item Fallbeschleunigung: $g \approx \SI{10}{\metre\per\second\squared}$
    \item Notiere bei Rechenaufgaben: Gegeben, Gesucht, Formel, Einsetzen, Ergebnis
\end{itemize}
}}
\end{center}

\vspace{0.5cm}

% ═══════════════════════════════════════════════════════════════
% AUFGABE 1: Grundwissen Kinematik (8 BE, AFB I)
% ═══════════════════════════════════════════════════════════════

\noindent\fbox{\parbox{\dimexpr\textwidth-2\fboxsep-2\fboxrule}{%
\textbf{Aufgabe 1: Grundwissen Kinematik}\hfill{\small\color{zeitgray}(ca.\ 9\,Min, AFB\,I)}\quad\textbf{8 BE}%
}}

\vspace{0.4cm}

\begin{enumerate}
\item[a)] \textbf{Nenne} die physikalische Größe, die man aus der Fläche unter einem $v$-$t$-Graphen ermitteln kann.

\hfill{\small\color{zeitgray}(1\,Min, AFB\,I)}\quad\textbf{(\_\_/1\,BE)}

\vspace{1.2cm}

\item[b)] \textbf{Gib an}, was die Steigung einer Geraden im $v$-$t$-Diagramm physikalisch bedeutet.

\hfill{\small\color{zeitgray}(1\,Min, AFB\,I)}\quad\textbf{(\_\_/1\,BE)}

\vspace{1.2cm}

\item[c)] \textbf{Rechne um:} \BEZEIT{2}{I}{2}

\vspace{0.3cm}
\hspace{1cm} 72 km/h = \underline{\hspace{3cm}} m/s

\vspace{0.5cm}
\hspace{1cm} 36 km/h = \underline{\hspace{3cm}} m/s

\end{enumerate}

\newpage

\begin{enumerate}
\item[d)] \textbf{Ergänze} die Tabelle. Trage die passende Graphenform ein. \BEZEIT{5}{I}{4}
\end{enumerate}

\vspace{0.5cm}

\begin{center}
\begin{tabular}{|p{6cm}|p{3cm}|p{5cm}|}
\hline
\textbf{Bewegungsart} & \textbf{Diagramm} & \textbf{Graphenform} \\
\hline
Gleichförmige Bewegung & $s$-$t$ & \\[0.8cm]
\hline
Gleichmäßig beschleunigte Bewegung & $s$-$t$ & \\[0.8cm]
\hline
Abbremsen (negative Beschleunigung) & $v$-$t$ & \\[0.8cm]
\hline
Stillstand & $v$-$t$ & \\[0.8cm]
\hline
\end{tabular}
\end{center}

\vspace{0.5cm}

\textit{Auswahl: Waagerechte Linie bei $v=0$ \quad|\quad Ansteigende Gerade \quad|\quad Fallende Gerade \quad|\quad Nach oben geöffnete Parabel}

\vspace{1cm}

% ═══════════════════════════════════════════════════════════════
% AUFGABE 2: Freier Fall -- Brunnen (6 BE, AFB I/II/III)
% ═══════════════════════════════════════════════════════════════

\noindent\fbox{\parbox{\dimexpr\textwidth-2\fboxsep-2\fboxrule}{%
\textbf{Aufgabe 2: Freier Fall -- Brunnenexperiment}\hfill{\small\color{zeitgray}(ca.\ 9\,Min, AFB\,I/II/III)}\quad\textbf{6 BE}%
}}

\vspace{0.4cm}

Um die Tiefe eines alten Brunnens zu bestimmen, lässt man einen Stein hineinfallen und misst die Zeit bis zum Aufprall.

\vspace{0.3cm}

\begin{enumerate}
\item[a)] Für den freien Fall gilt $h = \frac{1}{2}g \cdot t^2$.

\textbf{Stelle} diese Formel nach der Fallhöhe $h$ \textbf{um}, wenn die Fallzeit $t$ bekannt ist (die Formel ist bereits nach $h$ aufgelöst -- stelle stattdessen nach $t$ um). \BEZEIT{3}{I}{2}

\vspace{4cm}

\item[b)] \textbf{Berechne} die Tiefe des Brunnens, wenn der Stein $t = \SI{3}{\second}$ fällt.

\textit{Hinweis: Nutze die ursprüngliche Formel.} \BEZEIT{3}{II}{2}

\vspace{4.5cm}

\item[c)] In Wirklichkeit hört man den Aufprall etwas später, weil der Schall Zeit braucht, um nach oben zu gelangen. \textbf{Erkläre}, ob die tatsächliche Brunnentiefe größer oder kleiner ist als berechnet. \BEZEIT{3}{III}{2}

\vspace{3.5cm}

\end{enumerate}

\newpage

% ═══════════════════════════════════════════════════════════════
% AUFGABE 3: Bremsvorgang Fahrrad (8 BE, AFB I/II/III)
% ═══════════════════════════════════════════════════════════════

\noindent\fbox{\parbox{\dimexpr\textwidth-2\fboxsep-2\fboxrule}{%
\textbf{Aufgabe 3: Bremsvorgang eines Fahrrads}\hfill{\small\color{zeitgray}(ca.\ 11\,Min, AFB\,I/II/III)}\quad\textbf{8 BE}%
}}

\vspace{0.4cm}

Ein Fahrradfahrer fährt mit \SI{36}{\kilo\metre\per\hour} einen Hügel hinunter. Plötzlich läuft ein Hund auf die Straße und der Fahrer bremst mit $a = \SI{-5}{\metre\per\second\squared}$.

\vspace{0.3cm}

\begin{enumerate}
\item[a)] \textbf{Rechne} die Geschwindigkeit in m/s \textbf{um}. \BEZEIT{1}{I}{1}

\vspace{2cm}

\item[b)] \textbf{Berechne} die Bremszeit bis zum Stillstand. \BEZEIT{3}{II}{2}

\vspace{4cm}

\item[c)] \textbf{Berechne} den Bremsweg. \BEZEIT{3}{II}{2}

\vspace{4cm}

\item[d)] Der Hund ist \SI{12}{\metre} entfernt. \textbf{Beurteile}, ob der Fahrradfahrer rechtzeitig stoppen kann. \textbf{Begründe} Deine Antwort. \BEZEIT{4}{III}{3}

\vspace{4cm}

\end{enumerate}

\newpage

% ═══════════════════════════════════════════════════════════════
% AUFGABE 4: Straßenbahnfahrt (12 BE, AFB II/III)
% ═══════════════════════════════════════════════════════════════

\noindent\fbox{\parbox{\dimexpr\textwidth-2\fboxsep-2\fboxrule}{%
\textbf{Aufgabe 4: Straßenbahnfahrt}\hfill{\small\color{zeitgray}(ca.\ 25\,Min, AFB\,II/III)}\quad\textbf{12 BE}%
}}

\vspace{0.4cm}

Eine Straßenbahn fährt von einer Haltestelle zur nächsten. Die Fahrt verläuft in drei Phasen:

\begin{itemize}
\item \textbf{Phase 1 (Anfahren):} Die Bahn beschleunigt gleichmäßig in \SI{4}{\second} von $v = 0$ auf $v = \SI{12}{\metre\per\second}$.
\item \textbf{Phase 2 (Fahrt):} Die Bahn fährt \SI{8}{\second} lang mit konstanter Geschwindigkeit.
\item \textbf{Phase 3 (Bremsen):} Die Bahn bremst gleichmäßig in \SI{4}{\second} bis zum Stillstand.
\end{itemize}

\vspace{0.3cm}

\begin{enumerate}
\item[a)] \textbf{Zeichne} das $v$-$t$-Diagramm für die beschriebene Fahrt in das vorgegebene Koordinatensystem. Beschrifte die Achsen vollständig. \BEZEIT{6}{II}{3}

\vspace{0.3cm}

\begin{center}
\begin{tikzpicture}
\begin{axis}[
    width=14cm, height=5cm,
    xlabel={},
    ylabel={},
    xmin=0, xmax=18,
    ymin=0, ymax=15,
    grid=both,
    minor tick num=1,
    grid style={gray!20},
    major grid style={gray!35},
    axis lines=left,
    xtick={0,2,4,6,8,10,12,14,16,18},
    ytick={0,3,6,9,12,15},
    xticklabels={},
    yticklabels={},
    every axis label/.style={font=\footnotesize},
    tick label style={font=\footnotesize},
]
\end{axis}
\end{tikzpicture}
\end{center}

\vspace{0.5cm}

\item[b)] \textbf{Berechne} den gesamten zurückgelegten Weg der Straßenbahn. \BEZEIT{8}{II}{4}

\textit{Tipp: Der zurückgelegte Weg entspricht der \textbf{Fläche unter der Kurve} im $v$-$t$-Diagramm. Zerlege die Fläche in einfache geometrische Formen (Dreiecke, Rechtecke).}

\vspace{7cm}

\end{enumerate}

\newpage

\begin{enumerate}
\item[c)] \textbf{Zeichne} das zugehörige $s$-$t$-Diagramm. Beschrifte die $s$-Achse mit geeigneten Werten und achte auf die korrekte Kurvenform in jeder Phase. \BEZEIT{10}{III}{5}

\textit{Hinweis: Überlege, welche Kurvenform zu beschleunigter, gleichförmiger und verzögerter Bewegung gehört.}

\vspace{0.3cm}

\begin{center}
\begin{tikzpicture}
\begin{axis}[
    width=14cm, height=8cm,
    xlabel={$t$ in s},
    ylabel={},
    xmin=0, xmax=18,
    ymin=0, ymax=180,
    grid=both,
    minor tick num=1,
    grid style={gray!20},
    major grid style={gray!35},
    axis lines=left,
    xtick={0,2,4,6,8,10,12,14,16,18},
    ytick={0,20,40,60,80,100,120,140,160,180},
    yticklabels={},
    every axis label/.style={font=\footnotesize},
    tick label style={font=\footnotesize},
]
% Gestrichelte vertikale Phasengrenzen
\draw[gray!70, dashed, thick] (axis cs:4,0) -- (axis cs:4,180);
\draw[gray!70, dashed, thick] (axis cs:12,0) -- (axis cs:12,180);
\end{axis}
\end{tikzpicture}
\end{center}

\end{enumerate}

\newpage

% ═══════════════════════════════════════════════════════════════
% AUFGABE 5: Radfahrer überholt Jogger (8 BE, AFB I/II/III)
% ═══════════════════════════════════════════════════════════════

\noindent\fbox{\parbox{\dimexpr\textwidth-2\fboxsep-2\fboxrule}{%
\textbf{Aufgabe 5: Radfahrer überholt Jogger}\hfill{\small\color{zeitgray}(ca.\ 12\,Min, AFB\,I/II/III)}\quad\textbf{8 BE}%
}}

\vspace{0.4cm}

Im Park joggt eine Person mit konstant \SI{9}{\kilo\metre\per\hour}. Ein Radfahrer nähert sich von hinten mit konstant \SI{18}{\kilo\metre\per\hour}. Der Abstand zwischen beiden beträgt anfangs \SI{50}{\metre}.

\vspace{0.3cm}

\begin{enumerate}
\item[a)] \textbf{Rechne} beide Geschwindigkeiten in m/s \textbf{um}. \BEZEIT{1}{I}{1}

\vspace{2.5cm}

\item[b)] \textbf{Bestimme} die Relativgeschwindigkeit des Radfahrers bezüglich des Joggers.

\hfill{\small\color{zeitgray}(3\,Min, AFB\,II)}\quad\textbf{(\_\_/2\,BE)}

\vspace{3.2cm}

\item[c)] \textbf{Berechne}, nach welcher Zeit der Radfahrer den Jogger erreicht. \BEZEIT{3}{II}{2}

\vspace{3.5cm}

\item[d)] Nach \SI{100}{\metre} endet der gerade Parkweg und macht eine scharfe Kurve. \textbf{Analysiere}, ob der Radfahrer den Jogger vor der Kurve überholen kann. Für sicheres Überholen muss er den Jogger vollständig passiert haben (rechne mit \SI{10}{\metre} Sicherheitsabstand). \BEZEIT{5}{III}{3}

\vspace{5cm}

\end{enumerate}

\newpage

% ═══════════════════════════════════════════════════════════════
% AUFGABE 6: Kaffeefilter-Experiment (6 BE, AFB II/III)
% ═══════════════════════════════════════════════════════════════

\noindent\fbox{\parbox{\dimexpr\textwidth-2\fboxsep-2\fboxrule}{%
\textbf{Aufgabe 6: Kaffeefilter-Experiment}\hfill{\small\color{zeitgray}(ca.\ 10\,Min, AFB\,II/III)}\quad\textbf{6 BE}%
}}

\vspace{0.4cm}

Im Unterricht werden Kaffeefilter aus verschiedenen Höhen fallen gelassen. Dabei beobachtet man, dass die Filter nach kurzer Fallstrecke mit nahezu konstanter Geschwindigkeit sinken.

\vspace{0.3cm}

\begin{enumerate}
\item[a)] \textbf{Erkläre}, warum der Kaffeefilter direkt nach dem Loslassen zunächst beschleunigt wird.

\hfill{\small\color{zeitgray}(3\,Min, AFB\,II)}\quad\textbf{(\_\_/2\,BE)}

\vspace{3.7cm}

\item[b)] \textbf{Erkläre}, warum der Filter nach kurzer Zeit mit konstanter Geschwindigkeit fällt. Verwende die Begriffe \textit{Gewichtskraft} und \textit{Luftwiderstand}. \BEZEIT{4}{III}{2}

\vspace{5cm}

\item[c)] Ein Schüler faltet den Kaffeefilter zusammen und lässt ihn erneut fallen. \textbf{Begründe}, warum der gefaltete Filter schneller fällt als der ungefaltete. \BEZEIT{3}{III}{2}

\vspace{4cm}

\end{enumerate}

\vspace{1cm}

\begin{center}
\textit{--- Ende der Probeklausur ---}
\end{center}

\end{document}
